\documentclass[12pt]{article}

\newcommand{\I}{\mathbb{I}}
\newcommand{\Z}{\mathbb{Z}}
\newcommand{\N}{\mathbb{N}}
\newcommand{\Q}{\mathbb{Q}}
\newcommand{\R}{\mathbb{R}}
\newcommand{\D}{\mathbb{D}}


\usepackage{comment}
\usepackage{amsmath}
\usepackage{amssymb}  % for \nmid
\begin{document}


\centerline{\textbf{Gurkeerat Singh}}
\centerline{Homework 02, MORALLY Due Feb 17}

\newif{\ifshowsoln}
\showsolntrue % comment out to NOT show solution inside \ifshowsoln and \fi blocks.

\begin{enumerate}


\item
(40 points)
For this problem FIRST write the program and make conjectures THEN
look up whats true. 

\begin{enumerate}
\item
(0 points) 
Write a program that does the following:
For every $1\le n\le 900$:
\begin{itemize}
\item
Using the greedy method (Leo will tell you about that in class)
determine natural numbers $(x_1,x_2,\ldots,x_k)$ such that 
$x_1^3 + \cdots + x_k^3=n$. 
Using brute force (Leo will tell you about that in class)
determine natural numbers $(x_1,x_2,\ldots,x_L)$ such that 
$x_1^3 + \cdots + x_L^3=n$ with $L$ minimum. 
\item
Below I have the first 9 rows of the output. We use Gr-$k$ and Br-$L$ so the table fits on the page.
\end{itemize}

    

\hspace{-10em}\begin{tabular}{|l|l|l|l|l|}
\hline
 $n$ & Gr-$k$ & Greedy-$x_i$'s & Br-$L$ & Brute-$x_i$'s \cr
\hline
1   & 1         & $1=1^3$            & 1         & $1=1^3$       \cr
2   & 2         & $2=1^3+1$          & 2         & $2=1^3+1^3$       \cr
3   & 3         & $3=1^3+1^3+1^3$    & 3         & $3=1^3+1^3+1^3$   \cr
4   & 4         & $4=1^3+1^3+1^3+1^3$   & 4      & $4=1^3+1^3+1^3+1^3$   \cr
5   & 5         & $5=1^3+1^3+1^3+1^3+1^3$   & 5  & $5=1^3+1^3+1^3+1^3+1^3$   \cr
6   & 6         & $6=1^3+1^3+1^3+1^3+1^3+1^3$   & 6  & $6=1^3+1^3+1^3+1^3+1^3+1^3$   \cr
7   & 7         & $7=1^3+1^3+1^3+1^3+1^3+1^3+1^3$   & 7  & $7=1^3+1^3+1^3+1^3+1^3+1^3+1^3$   \cr
8   & 1         & $8=2^3$   & 1  & $8=2^3$   \cr
9   & 2         & $9=2^3+1^3$   & 2  & $9=2^3+1^3$   \cr
\hline
\end{tabular}

\bigskip

DO NOT hand in the program or the output. 

\vfill

\centerline{\bf GO TO NEXT PAGE FOR THE REST OF THIS PROBLEM}

\newpage

\item
(30 points)\begin{enumerate}
    \item (15 points)
Create a table that shows how many numbers $\leq 900$ require 1 cube, 2 cubes, 3 cubes, and so on for both the Greedy and Brute Force solutions. Below is a sample table (THIS DOES NOT CONTAIN THE CORRECT ANSWER):
\begin{center}
    \begin{tabular}{|c|c|c|}
     \hline  $\# $ of Cubes & Greedy & Brute Force  \\ \hline
        1 & 3 & 8 \\
        2 & 4 & 2 \\
        3 & 28 & 17 \\
        \vdots & \vdots &\vdots \\ \hline
    \end{tabular}
\end{center}
\begin{table}[h]
\centering
\begin{tabular}{|c|c|c|}
\hline
\# of Cubes & Greedy & Brute Force \\ 
\hline
1  & 9   & 9   \\ 
2  & 34  & 40  \\ 
3  & 70  & 113 \\ 
4  & 108 & 216 \\ 
5  & 130 & 258 \\ 
6  & 129 & 178 \\ 
7  & 121 & 69  \\ 
8  & 106 & 15  \\ 
9  & 91  & 2   \\ 
10 & 67  & 0   \\ 
11 & 29  & 0   \\ 
12 & 6   & 0   \\ 
\hline
\end{tabular}
\end{table}

\item (15 points) Based on this data make conjectures of the following forms:

\begin{enumerate}
\item 
Every $n$ is the sum of $\le XXX$ cubes. 
Write your conjecture in quantifiers. 
\begin{itemize}
    \item \boldmath{$\forall n \exists x_1 \exists x_2 \ldots \exists x_9 (n = (x_1)^3 + (x_2)^3 \ldots + (x_3)^3)$}
\end{itemize}
\item 
All but a finite number of $n$ is the sum of $\le XXX$ cubes. 
Write your conjecture in quantifiers.
\begin{itemize}
    \item \boldmath{$\exists N \forall n (n > N \rightarrow \forall n \exists x_1 \exists x_2 \ldots \exists x_9 (n = (x_1)^3 + (x_2)^3 \ldots + (x_3)^3))$}
\end{itemize}
\item
For large $n$, the Greedy algorithm gives the same answer as brute force for approximately XXX fraction of $\{1,...,n\}$. 
(Note that brute force gives the {\it correct} answer so we are wondering how often Greedy gives the
correct answer.)
\begin{itemize}
    \item \boldmath{$\lim_{x \to \infty} \frac{(k \leq n) s.t. P(k)}{n} = 0.43$}
    \item Assuming P(k) stands for "The greedy algorithm gives the same answer as brute force"
\end{itemize} 
\end{enumerate}
\end{enumerate}
\item
(10 points) 
Look on the web and/or use AI to determine what is known and what is conjectured
about these problems.
\begin{itemize}
    \item \textbf{Part A is known as Waring's Problem of Cubes, which states that every nonnegative integer can be written by a sum of most 9 cubes. Some numbers have lower answers but 9 is the known upper bound.} 
    \item \textbf{Part B is known as the Four Cubes Conjecture, which means that for large numbers n, n can be written by the sum of 4 cubes. the set of numbers that require more than 4 cubes is possibly finitely small. So past some threshold, every number is a sum of 4 cubes.}
    \item \textbf{Part C just measures the efficacy of greedy algorithms compared to brute force algorithm when it comes to finding the amount of cubes needed to generate some number n. There is no known conjecture or theory, and people's research on both these algorithms is mostly heuristic and experimental.}
    \item Data is sourced by ChatGPT
\end{itemize}
\end{enumerate}

\vfill

\centerline{\bf GO TO NEXT PAGE}

\newpage

\item
(30 points) 
In this problem I will give a property of a domain $\D$ and I want 
you to give me both: \begin{itemize}
    \item That property expressed as quantifiers.
    \item Either a domain $\D$, $\D \subseteq \R$, that has that property or the statement (without proof) that there is no such domain.
\end{itemize}

I give an EXAMPLE:

$\D$ has a least element

Expressed as quantifiers: $(\exists x\in \D)(\forall y\in \D)[x\le y]$.

Domain that works: $\N$ or you could write $\{0,1,2,\ldots\}$. 

\begin{enumerate}
\item
(10 points)
$\D$ is infinite and has both a least element and a greatest element.
\begin{itemize}
    \item \boldmath{$((\exists x)(\forall y)(x < y)) \wedge ((\exists w)(\forall z)(z \le w)) \wedge  (\forall a(\neg \forall b(b \le a) \rightarrow \exists c(a < c)))$}
    \item \boldmath{$\Q \bigcap  [0,1]$}
\end{itemize}

\item 
(10 points)
Every element $x\in \D$ has both a successor  element $y$ (so $x<y$ and there is nothing inbetween)
and also a predecessor element $w$ (so $w<x$ and there is nothing inbetween). 
\begin{itemize}
    \item \boldmath{$(\forall x \in \D)((\exists y)(x < y) \wedge (\neg \exists a)(x < a < y)) \wedge ((\exists w)(x > w) \wedge (\neg \exists b)(x > b > w))$}
    \item \textbf{$\Z$}
\end{itemize}

\item
(10 points)
$\D$ satisfies all the following properties
\begin{enumerate}
    \item $\D$ has a least element.
    \item $\D$ has a greatest element.
    \item Every element in $\D$ except the least has a predecessor.
    \item Every element in $\D$ except the greatest has a successor.
    \begin{itemize}
        \item \boldmath{$((\exists x)(\forall y)(x < y)) \wedge ((\exists w)(\forall z)(z \le w)) \wedge ( \forall a(a \neq x)(\exists c)(c < a) \wedge (\neg \exists d)(c < d < a) \wedge ( \forall e(e \neq w)(\exists f)(f < e) \wedge (\neg \exists g)(f < g < e))$}
        \item \textbf{\{1,2,3,4,5\}}
    \end{itemize}
\end{enumerate} (So the question here is to determine
if such a $\D$ exists, and either give me the $\D$ or state, without proof, there is no $\D$.)
\end{enumerate}


\ifshowsoln



\fi 

\vfill

\centerline{\bf GO TO NEXT PAGE} 

\newpage

\item
(30 points)
In this problem the only symbols you can use are 

\begin{itemize}
\item 
The usual logical ones: $\forall$, $\exists$ $\wedge$, $\vee$, $\neg$.
\item 
Variables that range over the domain: $x_1,x_2,x_3,\ldots$.
\item
These math symbols: $<$, $\le$, $>$, $\ge$, $=$, $\ne$.
\end{itemize}

For each of the following either give me the sentence I want OR
state (without proof) that there is no such sentence. 

\begin{enumerate}
\item
(15 points)
A sentence which is true with domain $\Z$ but false with domain $\Q$. 
\begin{itemize}
    \item \boldmath{$(\forall x_1)(\exists x_2)(x_1 \textless x_2 \wedge (\neg \exists x_3 (x_1 \textless x_3\wedge x_3 \textless x_2)))$}
\end{itemize}
\item 
(15 points)
($(0,1)$ is the set of all reals between 0 and 1 but not including 0 or 1.)
A sentence which is true with domain $(0,1)$ but false with domain $\Q$. 
(Warning: You can't use $(\exists x)[x^2=\frac{1}{2}]$ since this cannot be
stated just using $<$.) 
\begin{itemize}
    \item \textbf{No such sentence exists}
\end{itemize}
\end{enumerate}


\vfill

\centerline{\bf GO TO NEXT PAGE} 

\newpage
\item 
{\bf Honors HW 2} This is graded separately as Honors HW02.

Let $\I$ be the set of irrationals. As usual $\Q$ is the set of rationals.

\begin{enumerate}
\item 
(0 points, nothing to hand in)
Find someone in the class to play the DUP-SPOILER game $(\I,\Q,5)$.

\item
(100 points)
Determine which of the following statements is true and give an informal proof of it:

\begin{itemize}
\item 
There is a $k$ such that SPOILER wins $(\I,\Q,k)$.
The informal proof is a value $k\in\N$ and a strategy
for SPOILER to win $(\I,\Q,k)$. 
You DO NOT have to prove that
the strategy works, which is why I called it an {\it informal proof}.
BUT the strategy has to be well enough written so that I could carry
it out if asked to play the game.

\item
There is no $k$ such that SPOILER wins $(\I,\Q,k)$. 
The informal proof is to give, for all $k\in\N$, a strategy
for DUP to win $(\I,\Q,k)$. 
You DO NOT have to prove that
the strategy works, which is why I called it an {\it informal proof}.
BUT the strategy has to be well enough written so that I could carry
it out if asked to play the game.
\begin{itemize}
    \item \textbf{There is no such k such that SPOILER wins $(\I, \Q, k)$. We can define the irrational numbers as the set of numbers which are non-repeating and non-terminating. Similarly, the rational numbers can be defined by having repeating or terminating decimal expansions. Assuming both players play with the best-optimized strategy, the spoiler would try to trap the duplicator by trying to break the duplicator's set ordering. However, since both sets can include numbers infinitely small, it's not possible to break the ordering because continously smaller numbers in each set can be chosen that still fit natural order. In fact, both sets are structurally identical, making them isomorphic, and thus, there are no such amount of finite moves, including 5 moves, that could break the continuous property of both these sets.}
\end{itemize}
\end{itemize}


\end{enumerate}
\end{enumerate}

\end{document}

